%% @Author: KEMKA TAKENGNY Ulrich
%  @Date:   2024-08
%% @Class:  PFE IMIE-PARIS.

\documentclass[a4paper, oneside, 12pt, final]{extreport}
\usepackage{graphicx}

\parindent 0cm
\usepackage{makeidx}
\makeindex
%\usepackage[english]{babel}
\usepackage[lined,boxed,commentsnumbered, english, ruled,vlined,linesnumbered, french]{algorithm2e}
\usepackage{amsthm}
\newtheorem{theorem}{Theorem}[chapter]
\newtheorem{definition}{Definition}[chapter]
\newtheorem{exemple}{Example}[chapter]


%\usepackage[nottoc]{tocbibind}
%\addcontentsline{toc}{section}{References}

\providecommand{\keywords}[1]{\textbf{\textit{Mots clés---}} #1}
\providecommand{\keywordss}[1]{\textbf{\textit{Keywords---}} #1}

\usepackage{etoolbox}
%\makeatletter
%\patchcmd{\thebibliography}{%
%  \chapter*{\bibname}\@mkboth{\MakeUppercase\bibname}%{\MakeUppercase\bibname}}{%
%  \section{References}}{}{}
%\makeatother



\usepackage[nottoc]{tocbibind}

\textwidth 18cm
\textheight 24cm
\topmargin -1cm
\oddsidemargin -1cm

% set font encoding for PDFLaTeX or XeLaTeX
\usepackage{ifxetex}
\ifxetex
  \usepackage{fontspec}
\else
  \usepackage[T1]{fontenc}
  \usepackage[utf8]{inputenc}
  \usepackage{lmodern}
\fi


% Enable SageTeX to run SageMath code right inside this LaTeX file.
% documentation: http://mirrors.ctan.org/macros/latex/contrib/sagetex/sagetexpackage.pdf
%\usepackage{sagetex}


\newcommand{\reportTitle} {
  \textsc{Projet de Fin d'\'etudes}
}

\newcommand{\reportAuthor} {
  \textbf{Ulrich KEMKA TAKENGNY}
}

\newcommand{\reportSubject} {
Optimisation des Algorithmes de Deep Learning pour la Détection et la Reconnaissance en Temps Réel de Véhicules Militaires dans des Images et Vidéos
}



\newcommand{\studyDepartment} {%
  Entreprise d'accueil :
  
}

\newcommand{\IMIE} {%
\\ Ecole Supérieure d'informatique IMIE-Paris
}

%\newcommand{\codePFE} {% Reference
%  Code PFE%
%}


\newcommand{\AU} {
\centering \textbf{Année Académique 2023-2024}
}


\newcommand{\specialcell}[1]{%
  \begin{tabularx}{\textwidth}{@{}X@{}}#1\end{tabularx}%
}

%%%%%%%%%%%%%%%%%%%%%%%%%%%%%%%%%%%%%%%%%%%%%%%%%%%%%%%
% Add your own commands here
%%%%%%%%%%%%%%%%%%%%%%%%%%%%%%%%%%%%%%%%%%%%%%%%%%%%%%%
\newcommand{\MyCommand} {%
  Does nothing really%
}


% used in maketitle
\title{\reportSubject}
\author{\reportAuthor}

% Enable SageTeX to run SageMath code right inside this LaTeX file.
% documentation: http://mirrors.ctan.org/macros/latex/contrib/sagetex/sagetexpackage.pdf
%\usepackage{sagetex}

%\hypersetup{
%  pdftitle={\reportTitle~-~\reportSubject},%
%  pdfauthor={\reportAuthor},%
%  pdfsubject={\reportSubject},%
%  pdfkeywords={report} {internship} {pfe} {enis}
%}

\usepackage{graphics}
\usepackage{graphicx}


\usepackage[acronym,toc,section=chapter]{glossaries}
\makeglossaries

\newacronym{abc}{ABC}{A contrived acronym}
\newacronym{efg}{EFG}{Another acronym}
\newacronym{svm}{SVM}{Support Vector Machines}

\pagenumbering{roman} 

\usepackage[utf8]{inputenc}
%\usepackage[french]{babel}

\begin{document}
\thispagestyle{empty}
\begin{titlepage}
  \begin{center}


    %%%%%%%%%%%%%%%%%%%%%%%%%%%%%%%%%%%%%%%%%%%%%%%
    % THE HEADER
    %%%%%%%%%%%%%%%%%%%%%%%%%%%%%%%%%%%%%%%%%%%%%%%

    \includegraphics[scale=0.7]{./images/inr_logo_rouge.png}
    \hfill
    \includegraphics[scale=0.4]{./images/logonoir.png}

    \vspace{3cm}



    %%%%%%%%%%%%%%%%%%%%%%%%%%%%%%%%%%%%%%%%%%%%%%%
    % THE PAGE CONTENT
    %%%%%%%%%%%%%%%%%%%%%%%%%%%%%%%%%%%%%%%%%%%%%%%

    % \vspace{10pt} {%
    %   \renewcommand*{\familydefault}{\defaultFont}
    %   \fontsize{46pt}{46pt}\selectfont
    %   % MEMOIRE\\%
    %   %\reportTitle{}%\\\textsc{Report}\\%
    % }

    % \vspace{5pt}

    % \vspace{15pt}
    \textbf{\textit{Mémoire de fin d'études}}\\

    \vspace{5pt}
    {\textit{Réalisé par}}\\
    \vspace{10pt} {%
      \fontsize{14pt}{14pt}\selectfont
      {\bfseries\Large\sc \reportAuthor}\\
    }

    \vspace{10pt}

    \vspace{5pt} {
      \renewcommand*{\familydefault}{\defaultFont}
      \fontsize{27pt}{27pt}\selectfont%
      \rule{0.5\textwidth}{.4pt}\\
      \vspace{10pt}
      \reportSubject{}\\%
      \vspace{10pt}
      \rule{0.5\textwidth}{.4pt}
    }


    \vspace{30pt}
    {Diplôme préparé : \textbf{\large Manager de Solutions Digitales et Data}}\\
    %Soutenu le \dateSoutenance, devant la commission d'examen:\\
    \vspace{46pt}


    Encadrant Académique : \hfill Encadrant Professionnel :\\
    \vspace{10pt}
    \textbf{Mohamed EL AFRIT \hfill  Jonas RENAULT}


    %\vfill

    \vspace{40pt}
    % \textbf{\textit{Projet de Fin d'Etudes fait \`a}}\\
    % \vspace{5pt}
    % \studyDepartment
    % \centering
    % \includegraphics[scale=0.08]{./images/logonoir.png}
  \end{center}
  \vspace{40pt}
  \AU\\
\end{titlepage}

% ###############################
% # HELP COMMANDS               #
% ###############################
%
% -1 \part{part}
%  0 \chapter{chapter}
%  1 \section{section}
%  2 \subsection{subsection}
%  3 \subsubsection{subsubsection}
%  4 \paragraph{paragraph}
%  5 \subparagraph{subparagraph}


%%%%%%%%%%%%%%%%%%%%%%%%%%%%%%%%%%%%%%%%%%%%%%%%%%%%%%%
% Dédicace et Remerciements
%%%%%%%%%%%%%%%%%%%%%%%%%%%%%%%%%%%%%%%%%%%%%%%%%%%%%%%

\chapter*{Dedication}
%\chapter*{D\'edicace}
%\addcontentsline{toc}{chapter}{Dedication}
\thispagestyle{empty}
%
%For all they have endured to satisfy all my needs and wishes

\begin{center}
  {\it

    A ... pour son(leur) sacrifice et son(leur) soutien, \\
    en témoignage de mon infinie reconnaissance et mon profond attachement \\
    \vspace{1cm}
    A tous ceux qui me sont chers...

  }
\end{center}
%
%\nopagebreak{%
% And maybe a quote here
% \raggedright\hspace{5.75cm} To all of you,~\\
%\raggedright\hspace{7.75cm} I dedicate this work.
%  \raggedleft\normalfont\large\itshape{} \reportAuthor\par%
%}
%
%\cleardoublepage%

\chapter*{Acknowledgment}
%\chapter*{Remerciements}
%\addcontentsline{toc}{chapter}{Thanks}
\thispagestyle{empty}
%
%Au terme de ce travail (A l'issue de ce travail), je tiens à remercier M., Mme, Pr., Dr. pour sa disponibilité et ses conseils judicieux. \\

Je n'aurais jamais pu réaliser ce projet sans la précieuse aide et sans le soutien d'un grand nombre de personnes dont la générosité, la bonne humeur et l'intérêt manifestés à l'égard de mon PFE m'ont permis de progresser. \\

%En premier lieu, je tiens à remercier mon encadrant universitaire, \juryMemberFour{}, pour la confiance qu'il m'a accordée en acceptant d'encadrer ce travail, pour ses multiples conseils et pour toutes les heures qu'il a consacrées à diriger ce travail. \\ 

%Je souhaiterais exprimer ma gratitude à \juryMemberThree{}, pour m’avoir donné envie de réaliser un mémoire sur ... au sein de \og \studyDepartment \fg. Je le remercie également pour son accueil chaleureux à chaque fois que j'ai sollicité son aide, ainsi que pour ses multiples encouragements. J’ai été extrêmement sensible à ses qualités humaines d'écoute et de compréhension tout au long de ce travail de mémoire. \\

%J'aimerais également dire à \juryPresident{} à quel point je suis honorée pour avoir accepté de présider ce jury de PFE. Je suis infiniment gré à  \juryMemberOne{} de s’être rendu disponible et d’avoir accepté la fonction de rapporteur. De même, je suis particulièrement reconnaissant(e) à (\juryMemberTwo{} de l'intérêt qu'il/elle a manifesté à l'égard de ce projet en s'engageant à être rapporteur. \\

Ma reconnaissance va à ceux qui ont plus particulièrement assuré le soutien affectif de ce travail : ma famille ainsi que mes amis. Mes parents...


\markboth{\MakeUppercase{Abstract}}{}
\addcontentsline{toc}{chapter}{Abstract}
\sloppy

\textbf{Abstract Provisoire}\\

This thesis focuses on the optimization of deep learning algorithms for the real-time detection and recognition of military vehicles in images and videos.
The primary objective is to study the state-of-the-art methods and propose improvements based on the integration of advanced techniques such as generative models and data augmentation.
The ADOMVI project, to which this work has contributed, aims to optimize the performance of detection models like YoloV8, considering the specific constraints of the military domain, including data scarcity and varying visibility conditions.
The results demonstrate a significant improvement in detection and recognition performance while highlighting persistent challenges related to the robustness of models in real-world conditions.
This thesis concludes with the necessity to continue research to refine these algorithms and adapt technological solutions to the needs of defense operations.



%%%%%%%%%%%%%%%%%%%%%%%%%%%%%%%%%%%%%%%%%%%%%%%%%%%%%%%
% Divers chapitres
%%%%%%%%%%%%%%%%%%%%%%%%%%%%%%%%%%%%%%%%%%%%%%%%%%%%%%%

\tableofcontents
%\addcontentsline{toc}{chapter}{\contentsname}

\listoffigures
%\addcontentsline{toc}{chapter}{Liste des Figures}
\listoftables
%\addcontentsline{toc}{chapter}{Liste des Tableaux}
\listofalgorithms
\addcontentsline{toc}{chapter}{Liste des algorithmes}
\cleardoublepage

\newpage
\pagenumbering{arabic}
\chapter*{Introduction}
\label{chap:general_intorduction}



\markboth{\MakeUppercase{Introduction}}{}%
\addcontentsline{toc}{chapter}{Introduction}%

%Welcome to \Ac{ITBS}. ~\\
%Again, welcome to \Ac{ITBS}. ~\\
%Your introduction goes here. ~\\

Voici une référence à l'image de la Figure \ref{fig:test} page \pageref{fig:test} et une autre vers la partie \ref{chap:2} page \pageref{chap:2}.
On peut citer un livre\, \cite{caillois1} et on précise les détails à la fin du rapport dans la partie références.
Voici une note\,\footnote{Texte de bas de page} de bas de page\footnote{J'ai bien dit bas de page}. Nous pouvons également citer l'Algorithme , la Définition \ref{def1}, le Théorème \ref{theo1} ou l'Exemple \ref{exo1}...\\

Le document est détaillé comme suit : le chapitre \ref{chap:chapterone} introduit le cadre général de ce travail. Il s'agit de présenter l'entreprise d'accueil et de détailler la problématique. Le chapitre \ref{chap:2} introduit les données ainsi que les modèles choisies.\\


Dans le cadre de notre formation et en tant qu’alternant, nous sommes partagés entre les
enseignements à l’école et les enseignements en entreprise, ceci pour faciliter notre immersion
dans le monde du travail. Pour mon alternance, j’ai intégré le centre de recherche Inria, dans le département Mission Défense et Sécurité en tant que Développeur Web
Au cours de cette expérience professionnelle, j’ai participé au développement d'une application web permettant contenant des scénarios permettant de simuler l'exploitation du renseignement d'intérêt militaire,
au déploiement d'un mail serveur simulant un scénario de sécurité économique dans un environnement d'entreprise (startup, équipe projet recherche) et
à l'optimisation des Algorithmes de Deep Learning pour la Détection et la Reconnaissance en Temps Réel de Véhicules  Militaires dans des Images et Vidéos
Ces projets dont l’objectif d’une part à faire appréhender aux académiques et aux entreprises les problèmes
concrets rencontrés par les opérationnels afin d’y proposer des solutions communes, et
d’autre part, de permettre d’expérimenter les solutions sur la base de procédures de tests
« opérationnels » constitueront la pierre angulaire notre rapport.



\chapter{Données étudiés}%
\label{chap:chapterone}
%%%%%%%%%%%%%%%%%%%%%%%%%%%%
% SECTION                  %
%%%%%%%%%%%%%%%%%%%%%%%%%%%%
\section{Présentation du centre de recherche}
\label{chap:sectionone}

\subsection{Inria : métiers et chiffres clés.}

L’Inria est l’institut national de recherche en sciences et technologies du numérique, crée en 1967, il dispose de 11 centres et plus de 20 antennes et emploie 2600 personnes. La recherche de rang mondial, l’innovation technologique et le risque entrepreneurial constituent son ADN.
Au sein de 215 équipes-projets en général communes avec des partenaires académiques, plus de 3 900 chercheurs et ingénieurs y explorent des voies nouvelles. 900 personnels d’appui à la recherche et à l’innovation contribuent à faire émerger et grandir des projets scientifiques ou entrepreneuriaux qui impactent le monde.

L’institut fait appel à de nombreux talents dans plus d’une quarantaine de métiers différents, travaille avec de nombreuses entreprises et a accompagné la création de plus de 180 start-ups.
Inria soutient la diversité des voies de l’innovation : de l’édition open source de logiciels à la création de startups technologiques (Deeptech).


\subsection{Le service de l’apprenti : Mission Défense et Sécurité}

Le renforcement des partenariats avec la sphère Sécurité/Défense de l’État est une priorité stratégique de l’Inria. C’est de ce contexte qu’est né l’institut. Créé en mars 2020 et dirigé par Frédérique Segond, la Mission Défense et Sécurité a pour objectif le soutient des politiques gouvernementales qui visent la souveraineté et l’autonomie stratégique numérique de l’Etat français, voire européen. Elle fédère tous les projets sécurité défense d’Inria bientôt de toute la France.
L’équipe en pleine croissance est actuellement composée de seize personnes ayant chacun un rôle bien définit avec le soutien des intervenants externes à la mission.

\section{Description du projet IntelLab}
\subsection{Contexte}

Les ambitions de la mission sont aujourd’hui implémentées à travers deux structures

\begin{itemize}\addtolength{\itemsep}{-0.35\baselineskip}
  \item IntelLab, un environnement de \textbf{simulation, formation, et d’expérimentation}, il vise d’une part à faire appréhender aux académiques et aux entreprises les problèmes concrets rencontrés par les opérationnels afin d’y proposer des solutions communes, et d’autre part, de permettre d’expérimenter les solutions sur la base de procédures de tests « opérationnels ».
  \item Un centre d’excellence dédié au domaine de la Sécurité / Défense afin de faciliter le développement et le transfert à court, moyen et long terme de technologies issues de la Recherche.
\end{itemize}

\noindent Entre autres, La plateforme IntelLab permet de jouer différents types de scénario :

\begin{itemize}\addtolength{\itemsep}{-0.35\baselineskip}%
  \item Des scénarios simulant l'exploitation du renseignement d'intérêt militaire ;
  \item Un scénario de Lutte Informatique d'influence (L2I) ;
  \item Des scénarios de détection de signaux faibles précurseurs d'actions d'ingérence.
\end{itemize}

\subsection{Principales activités}


Sous la supervision du responsable informatique, j'ai développé divers outils pédagogiques et technologiques visant à faciliter la prise en main des applications et à enrichir la compréhension des scénarios d'exercice de simulation. Mon travail a contribué à l'élaboration et à l'optimisation de nouvelles technologies au sein de l'équipe.

J'ai principalement participé au développement et à l'évolution de la plateforme IntelLab, un projet clé permettant de simuler des scénarios de sécurité variés.
Nous avons travaillé en mode projet avec la méthode Agile, ce qui nous a permis d'assurer une gestion efficace et un suivi optimal des projets.

\noindent Mes contributions se sont réparties sur trois projets distincts :

\begin{itemize}\addtolength{\itemsep}{-0.35\baselineskip}%
  \item \textbf{Développement de l'application web IntelLab }: J'ai pris en charge le développement des interfaces web Frontend et Backend de la plateforme, intégrant les dernières technologies pour garantir une performance et une expérience utilisateur optimales.
  \item \textbf{Mise en place d'un serveur de messagerie} : En parallèle, j'ai contribué au déploiement d'une plateforme dédiée à la simulation de scénarios de sécurité économique, centrée sur la gestion d'un serveur de messagerie. Cette plateforme est utilisée dans des exercices de simulation reproduisant des environnements d'entreprises tels que des startups ou des équipes de recherche.
  \item \textbf{Détection et reconnaissance, proche du temps réel, de véhicules militaires sur des images et vidéos} : J'ai participé à un projet visant à développer et affiner des algorithmes de Deep Learning pour la détection et la reconnaissance en temps quasi réel de véhicules militaires dans des images et vidéos.
\end{itemize}

\noindent Dans le cadre de ces projets, mes responsabilités incluent :

\begin{itemize}\addtolength{\itemsep}{-0.35\baselineskip}%
  \item \textbf{Développement web Frontend et Backend} : Conception et implémentation des fonctionnalités sur les deux volets de l'application.
  \item \textbf{Développement d'algorithmes génératifs} : Création et intégration d'algorithmes destinés à enrichir les jeux de données pour améliorer les performances des modèles.
  \item \textbf{Rédaction des tests} : Rédaction et exécution de tests pour garantir la fiabilité et la robustesse des solutions développées.
  \item \textbf{Mise en production} : Déploiement des solutions dans l'environnement de production, assurant leur bon fonctionnement et leur maintenance.
  \item \textbf{Maintenance des applications et de leurs infrastructures} : Assurer la stabilité et la disponibilité continue des applications en identifiant et résolvant les problèmes techniques rapidement.
\end{itemize}



\chapter{Modèles utilisés et Applications}
\label{chap:2}

  

\begin{itemize}
	\item The individual \index{Entries}{entries} are indicated with a black dot, a so-called bullet.
	\item The text in the entries may be of any length.
\end{itemize}

\begin{theorem}\label{theo1}
Soit $n$ un entier naturel. Si $n$ est premier alors il n'est divisible que par 1 et par lui-même.
\end{theorem}

\begin{proof}
Here is my proof.
\end{proof}

\begin{definition}\label{def1}
Soit $A$ une courbe...
\end{definition}

Ici, il s'agit de l'utilisation de TB %\nomenclature[TB]{TB}{Très Bien} qui consiste à parler Très Bien. 
\gls{abc} et \gls{efg} sont des acronyms et des abbréviations... La méthode \gls{svm} est également couramment utilisée.

\begin{exemple}\label{exo1}
On considère le cas particulier... 
\end{exemple}


\chapter*{Conclusion et Perspectives}
\label{chap:conclusion}
\markboth{\MakeUppercase{Conclusion}}{}%
\addcontentsline{toc}{chapter}{Conclusion}
And a very interesting conclusion here\@. ~\\
Lorem ipsum dolor sit amet, consectetur adipisicing elit, sed do eiusmod
tempor incididunt ut labore et dolore magna aliqua. Ut enim ad minim veniam,
quis nostrud exercitation ullamco laboris nisi ut aliquip ex ea commodo
consequat.

\newpage
\appendix
\addcontentsline{toc}{chapter}{Annexes}
%\markboth{\MakeUppercase{Annexe}}{}

\chapter{Code R pour résoudre la problématique}
\label{chap:appendix}


\section{Pré-traitement des données}
\section{Code R pour les modèles}

An appedix if you need it.

\begin{verbatim}
 Insérer ici le code !
 \end{verbatim}

\section{Librairies utilisées}

Lorem ipsum dolor sit amet, consectetur adipisicing elit, sed do eiusmod
tempor incididunt ut labore et dolore magna aliqua. Ut enim ad minim veniam,
quis nostrud exercitation ullamco laboris nisi ut aliquip ex ea commodo.


%%%%%%%%%%%%%%%%%%%%%%%%%%%%%%%%%%%%%%%%%%%%%%%%%%%%
% Don't touch this, it is auto generated
%%%%%%%%%%%%%%%%%%%%%%%%%%%%%%%%%%%%%%%%%%%%%%%%%%%%
\nocite{*}

%\phantomsection{}
%\addcontentsline{toc}{chapter}{Webography}
%\printbibliography[title={Webography},type=online]

%\phantomsection{}
%\addcontentsline{toc}{chapter}{Bibliography}
%\printbibliography[title={Bibliography},nottype=online]

%\printbibheading %exemple de bibliographie divisée en sections. Pour ajouter des oeuvres non citées,utiliser \nocite

%\printbibliography[keyword=pratique,heading=subbibliography,title={Théories littéraires dans les jeux vidéo}]
%\printbibliography[keyword=litteraire,heading=subbibliography,title={Narratologie et structuralisme}]

%\printbibliography[keyword=jeu,heading=subbibliography,title={\emph{Games studies}}]

\bibliographystyle{apalike}
%\bibliographystyle{plain}

\bibliography{Biblio.bib}

\cleardoublepage%

\addtocontents{toc}{\protect\setcounter{tocdepth}{3}}

\printglossaries
\printindex

\markboth{\MakeUppercase{Abstract}}{}
\addcontentsline{toc}{chapter}{Abstract}
\sloppy

\textbf{Abstract Provisoire}\\

This thesis focuses on the optimization of deep learning algorithms for the real-time detection and recognition of military vehicles in images and videos.
The primary objective is to study the state-of-the-art methods and propose improvements based on the integration of advanced techniques such as generative models and data augmentation.
The ADOMVI project, to which this work has contributed, aims to optimize the performance of detection models like YoloV8, considering the specific constraints of the military domain, including data scarcity and varying visibility conditions.
The results demonstrate a significant improvement in detection and recognition performance while highlighting persistent challenges related to the robustness of models in real-world conditions.
This thesis concludes with the necessity to continue research to refine these algorithms and adapt technological solutions to the needs of defense operations.





\end{document}
